\begin{center}{\Large\bfseries Response to the Independent Referees}\par\vspace{6pt}
Finite-Catalog Renewal Design by Selected-Boundary Resource Paths\par
R53, September 26, 2026\end{center}

\section*{Version and Principal Changes}
We thank both referees for identifying the substantive distinction between a correct certificate and a satisfactory structural algorithm. This response addresses the first and second independent reports on R49, with particular attention to the second report's explicit dyadic-count error and endogenous-catalog objection. The latest available report concerns R49, not R52 or the present revision. We therefore build on the intervening R52 manuscript rather than discarding its stronger algorithm and experiments.

The current article retains the original finite-catalog joint problem: individual targets remain endogenous, the accepted root obligation is met exactly, every expected cap and realization ceiling is enforced, and the installed command budget is unchanged. The arbitrary-eligibility result is a polynomial additive scheme, while common-prefix optimization remains exact. R53 adds two structural results to that retained foundation. Proposition~\ref{prop:potential53} identifies an eligibility-independent capacity potential for every arc. Theorem~\ref{thm:screen53} gives a uniform command-deletion envelope and simultaneous safe screening; Corollary~\ref{cor:transfer53} then transfers a reduced-catalog global certificate back to the original problem. These are proved and tested additions, not relabelings of an unchanged incumbent.

\section{Incorrect Dyadic Cover Count: Second Report 3 and 13.1}
We agree that the old leaf bound is false. With one free coordinate and width-to-tolerance ratio three, midpoint refinement needs four leaves, not three. The retained corrected statement in the electronic companion defines the weighted initial widths and uses the exact product
\[
 \prod_{g\ne g_0}2^{\max\{0,\lceil\log_2(w_g^0/\eta)\rceil\}}.
\]
It then gives the valid factor-two coarse bound and the binary-tree node count. The ratio-three example is explicit. Twenty exact combinatorial regressions include zero widths and anisotropic widths. The former count is preserved only in immutable historical sources, with its correction identified; it is not used by the current theorem, abstract, or algorithm. The current scalar-resource proof does not depend on that box-cover count.

\section{Catalog Instability and Unused Commands: Second Report 4, 13.2, and 13.3}
The referees are right that a full-catalog eligibility count is not an intrinsic dimension of an optimal selected book. The article now calls it a sufficient catalog-dependent parameter of the older class-mean method. The selected-boundary path of Theorem~\ref{thm:path52} groups histories at the commands actually selected, not at every unselected candidate. Retained path payoffs are unchanged when unrelated catalog points are inserted.

The catalog-stability proposition states that one insertion increases the full-catalog class count by at most one. Threshold perturbations that cross no candidate leave the canonical problem unchanged; an explicit crossing example shows why a continuity claim without a margin would be wrong. Endpoint certificates provide a robust bracket for uncertain ceilings. Proposition~\ref{prop:potential53} strengthens the analysis: the capacity of each resource arc is a difference of the same cap potential, independent of eligibility as long as the stipulated cap--ceiling relation holds. Eligibility changes its payoff, not this conserved capacity.

R53 addresses the specific unused-candidate objection constructively. The anchor envelope bounds the contribution of a non-anchor command in every possible selected neighboring context. Its opening payment can then certify exclusion without first solving the enlarged problem. The bound is invariant to other non-anchor insertions that retain the minimum catalog command and original primitives. All certified commands may be deleted simultaneously at unchanged individual targets. The accompanying transfer certificate checks those exclusions and the entire reduced global proof. An unchanged winning book alone is explicitly insufficient. Equality-charge deletion preserves an optimum, whereas strict inequality excludes the command from every optimum; the distinction is demonstrated by an exact example.

\section{A Stronger Variable-Eligibility Algorithm: Second Report 7 and 13.4}
We respond through the stronger-algorithm route requested by the report, not by assuming the old exponential dependence to be necessary. Theorem~\ref{thm:approx52} has a scalar resource grid and polynomial arithmetic and bit complexity in the explicitly stated parameters, with no exponent depending on the number of eligibility classes. Its repair step restores the original promise on the same selected path. It permits rational nonnegative charges, heterogeneous caps and quadratic curvatures including zero, interior promises, and arbitrary prefix eligibility.

The complexity statement separates numerical inverse normalized accuracy from binary accuracy encoding. It does not assert an exact polynomial algorithm for unrestricted budget or a relative-error guarantee near zero values. The exact common-prefix theorem is retained as a stronger result in its regime, not weakened to an approximation. R53 screening computes all command envelopes in linear history--catalog arithmetic time before this algorithm. It neither introduces another target grid nor expands the command budget. The independent screening checker intentionally uses a redundant polynomial context enumeration, and its cost is reported separately.

\section{Pooling Proof Details and Rational Encoding: Second Report 2}
The weighted transfer is displayed explicitly, including the recipient increment scaled by the probability ratio. The exact pooling proposition states equality and continuity at the aggregate terminal-capacity threshold. Its proof explains why a mean above that threshold forces the water level above the last eligible command, and writes the resulting residual-service identity. The companion includes both endpoint identities needed for the group-box service interval and all price-sign cases. The mathematical identities allow real nonnegative charges; every exact rational bit-complexity assertion requires rationally encoded charges. These details remain in the current manuscript and companion.

\section{Nontrivial Eligibility and Hard-Regime Evidence: Second Report 5, 6, 9, 10, 13.5, and 13.6}
The twelve retained large common-prefix cases use distinct ceilings strictly inside an interior catalog gap and exclude a nonempty suffix. They extend through 1,024 histories, vary catalog size, budget, and promise, and include heterogeneous weights, caps, charges, and zero curvatures. The original all-levels-eligible family is retained only as historical arithmetic evidence. A controlled eligibility sweep, threshold perturbations, and free and charged refinements are separately reported.

The retained study contains all 44 prescribed requests, including the six difficult models at two accuracies. The fine-tolerance resource intervals attain the requested width. Forty-two subdivision comparisons and twelve numerical mixed-integer envelope solves expose the competing costs. The article reports relative and absolute widths, certificate storage and bit lengths, checking costs, and a compact mixed-integer diagnostic table. Every historical unresolved request remains unchanged. Fixed-node subdivision comparisons are not mislabeled as equal-time or equal-memory comparisons; the numerical comparison has its declared measured time allowance. Numerical solver status is not used as a rational certificate.

The new twelve-case screening study tests cheaper inserted commands than the old gross exclusion rule can remove. It preserves every exhaustive optimum and reduces the resource tables even when the full-catalog class count rises to nine. All inputs, global intervals, winning books, timings, and certificate sizes are retained. Transfer certificates sometimes exceed the full certificate in small cases because they also retain the original instance and exclusion record. We report this overhead rather than select only favorable compression cases. This controlled study at width $1/10$ supplements, and does not replace, the earlier $1/1000$ study.

The new random regression covers 120 models, 2,332 capacity identities, 542 complete neighboring-context maxima, and 129,448 pointwise deletion comparisons. Every before--after exhaustive optimum agrees; 149 commands are removed. Twenty original-instance transfers pass, and fourteen altered certificates fail. These tests supplement the proofs. One initial test expected all four base candidates to remain; the mathematically valid removal of a wholly ineligible zero-charge endpoint required correcting that expectation. The initial diagnostic is preserved, and no experimental input or exclusion rule was changed to obtain the reported result.

\section{A Second Operations Research Model: Second Report 8 and 13.8}
The compatible-production model is stated independently of renewal histories and terminal lotteries. A compatible path installs modules whose divisible operating capacities jointly serve a fixed order, with concave returns and installation payments. Capacity-conservative configurations satisfy the same repair condition, yielding an original-order and original-module-budget guarantee. The article explicitly distinguishes this parallel capacity assignment from serial-flow bottlenecks. The new capacity-potential proposition now shows the exact structural connection to the renewal arcs, rather than offering only a verbal change of application labels.

The application is a mathematical extension, not a claim of industrial calibration. All computational inputs remain synthetic. The theorem and experiments establish the asserted optimization and certification properties; they do not measure field thresholds or the economic value of a chosen tolerance. The source and response maintain that distinction without dropping the joint-design contribution.

\section{Focused Article and Preservation: Second Report 12 and 13.7}
The journal-facing article centers on the canonical response, selected-boundary identity, conserved capacity, exact common-prefix algorithm, arbitrary-eligibility additive scheme, refinement-safe screening, and the associated evidence. The companion supplies the retained class-box proof, corrected dyadic analysis, certificate mathematics, and detailed diagnostics. Continuous-design, institutional, coarsening, and historical results remain intact in their original repository files and supplementary readers; they are not silently deleted or turned into additional prerequisites for this article.

The revision uses an anonymous title page, a single text-only abstract below 200 words, an equation-free introduction, author--year references, numbered tables following the references, and a code/data statement after the body. The build record reports actual page counts and rejects unresolved references, duplicate labels, and overfull boxes. Prior root readers are preserved byte for byte. A full parent-tree audit distinguishes preserved artifacts from the small set of replaced current-reader entry points. Successful compilation is not presented as an editorial decision.

\section*{Disposition of the Remaining Minor Comments}
The title identifies finite-catalog design. The abstract separates exact common-prefix optimization from the general additive result. The command alphabet is distinguished from target tables, encoder logic, numerical precision, and certificate storage in the model itself. Rational charges, strict-prefix geometry, relative errors, initial dyadic widths, and the ratio-three counterexample are explicit. Positive-width intervals are not called exact solutions; exhaustive attainment and zero-width certificates are distinguished. Post-optimization pooling preserves every individual witness and does not prove optimality over omitted books. A global certificate is reused after refinement only with the new checked exclusion argument or a fresh proof over the enlarged path set. The current README, source hashes, and preservation map identify exactly which version and evidence were used.
